\section{An exact finite-threshold recurrence}\label{sec:recurrence}

Select one run position in \eqref{eq:input}, and keep its generator,
its sign \(\varepsilon\), and every other run fixed. Only its absolute
exponent \(m\ge1\) varies. Define
\begin{equation}\label{eq:context-span}
a=-\sum_{\substack{j\ne *\\e_j<0}}|e_j|,\qquad
c=\sum_{\substack{j\ne *\\e_j>0}}|e_j|,\qquad
W=c-a,\qquad m_0=W+2.
\end{equation}
The letter \(c\) here is a degree endpoint, not a circle count or the
number \(b\) of braid strands. Empty sums are zero.

\subsection{The fixed context and its repeated differential}

For each context degree \(q\), let \(I^q\) be the direct sum of reduced
resolution spaces when the distinguished box contains \(I\), and let
\(E^q\) be the corresponding sum with \(E\) in that box. Both vanish
outside \([a,c]\). All other run differentials define \(D_I\) and \(D_E\).
The chosen saddle maps between the two sectors in the direction determined
by \(\varepsilon\). Write
\begin{equation}\label{eq:context-space}
V=\bigoplus_{q=a}^{c}E^q,\quad M=\dim V,\quad
L=X_p+X_q,\quad F=D_E+L,\quad
r=\rank F,\quad \eta=M-2r.
\end{equation}
In \(X_p+X_q\), the subscripts designate the two exterior arc points,
independently of the context degree index.
The operator \(D_E\) raises context degree, while \(L\) preserves it.
Consequently \(F\) is initially an endomorphism of the \emph{ungraded}
space \(V\).

\begin{lemma}\label{lem:square-zero}
The operator \(F\) satisfies \(F^2=0\), and
\(\eta=\dim H(V,F)\ge0\).
\end{lemma}

\begin{proof}
The context is a chain complex, so \(D_E^2=0\). Each dot squares to zero,
and distinct dots commute; hence \(L^2=0\) in characteristic two.
Multiplication and comultiplication commute with a dot on an unchanged
arc:
\[
m(X\otimes1)=Xm=m(1\otimes X),\qquad
(X\otimes1)\Delta=\Delta X=(1\otimes X)\Delta.
\]
These are direct identities in \(A\), including on the marked-\(X\)
subcomplex. Thus \(D_EL=LD_E\), and the mixed terms in \(F^2\) cancel.
Finally, \(\operatorname{im}F\subseteq\ker F\); rank--nullity gives
\(\dim\ker F-\dim\operatorname{im}F=M-2r\).
\end{proof}

The identity \(F^2=0\) is stronger than a numerical observation about
consecutive ranks. Its use is specific to the characteristic-two normal
form just specified. An integral or odd-characteristic version needs its
own signs, gradings, and differential analysis.

\subsection{Where complete context groups occur}

For a positive distinguished run, the summand \(E^q\) lies in total degree
\(h=q+k\), with \(1\le k\le m\). Therefore
\begin{equation}\label{eq:positive-central}
C_m^h=V\quad\text{throughout}\quad c+1\le h\le a+m.
\end{equation}
Indeed, these inequalities are exactly the conditions that every
\(q\in[a,c]\) has \(1\le h-q\le m\). The \(I\) sector is absent in
this interval because \(h>c\).
Between two adjacent groups satisfying \eqref{eq:positive-central},
the other-run maps give \(D_E\) and the distinguished internal maps give
\(L\). Under the canonical identification by context state and circle
label, the differential is exactly \(F\).

For a negative distinguished run, \(h=q-k\) instead. Complete groups and
their internal differentials occur in
\begin{equation}\label{eq:negative-central}
c-m\le h\le a-1,\qquad
\partial_h=F\quad(c-m\le h\le a-2).
\end{equation}
Here the \(I\) sector is absent because \(h<a\).
These descriptions keep the mixed signs in the context; no assumption
that the braid is positive, alternating, or homogeneous has entered.

At \(m=m_0\), there are exactly two consecutive complete groups. In the
positive case they are in degrees \(c+1,c+2\); in the negative case they
are in degrees \(a-2,a-1\). The single map between them is \(F\).
The two-group threshold is deliberate. One complete group alone does not
expose a complete central differential, and its Betti number can still
depend on both boundary regions.

\subsection{The two signed profile formulas}

Let
\[
P_m(z)=\sum_h\dim_{\F}H^h(C_m)\,z^h
\]
be the reduced homological Laurent polynomial, with quantum grading
forgotten and the macro grading fixed as in \cref{prop:normal}.

\begin{theorem}[One-reference recurrence]\label{thm:recurrence}
Set \(m_0=W+2\), \(\delta=m-m_0\), and compute \(P_{m_0}\).
For a positive distinguished run, split it into
\[
P_{\rm low}=\sum_{h\le c+1}\beta_h(m_0)z^h,\qquad
P_{\rm high}=\sum_{h\ge c+2}\beta_h(m_0)z^h.
\]
Then for every \(m\ge m_0\),
\begin{equation}\label{eq:positive-profile}
\boxed{\displaystyle
P_m(z)=P_{\rm low}(z)+z^\delta P_{\rm high}(z)
       +\eta\sum_{h=c+2}^{a+m-1}z^h.}
\end{equation}
For a negative distinguished run, put
\[
Q_{\rm low}=\sum_{h\le a-2}\beta_h(m_0)z^h,\qquad
Q_{\rm high}=\sum_{h\ge a-1}\beta_h(m_0)z^h.
\]
Then
\begin{equation}\label{eq:negative-profile}
\boxed{\displaystyle
P_m(z)=z^{-\delta}Q_{\rm low}(z)+Q_{\rm high}(z)
       +\eta\sum_{h=c-m+1}^{a-2}z^h.}
\end{equation}
An interval sum with lower endpoint above its upper endpoint is empty.
In either case the total reduced rank obeys
\begin{equation}\label{eq:rank-recurrence}
R(m)=R(m_0)+(m-m_0)\eta.
\end{equation}
The statement permits links and any fixed exterior reduction point.
\end{theorem}

\begin{proof}
First suppose the distinguished run is positive. Near the lower end,
the condition \(k\le m\) is inactive once \(m\ge m_0\).
Every chain group in degrees at most \(c+1\), its incoming map, and its
outgoing map are therefore the same as at \(m_0\).
At the top, writing \(k'=k-\delta\) identifies all groups in degrees at
least \(c+2+\delta\), together with both incident maps, with the reference
upper region translated by \(\delta\). This can also be read directly
from \(1\le h-q\le m\): near that end the lower bound is inactive.
The saddle into the first \(E\) sector remains in the fixed lower region.

The degrees strictly between these two regions are
\(c+2,\ldots,a+m-1\). Each has chain space \(V\), incoming map \(F\),
and outgoing map \(F\). Its Betti number is \(M-r-r=\eta\).
There are \(m-m_0\) such degrees. Applying \eqref{eq:betti} in the fixed
and translated boundary regions proves \eqref{eq:positive-profile}.

For a negative run, use \(1\le q-h\le m\). The lower region now
translates downward by \(\delta\), while the upper region, including
the saddle into \(I\), remains fixed. Its separating complete groups
are \eqref{eq:negative-central}. The inserted degrees are precisely
\(c-m+1,\ldots,a-2\), again with both maps \(F\).
This proves \eqref{eq:negative-profile}. Evaluating at \(z=1\) gives
\eqref{eq:rank-recurrence}.
\end{proof}

This proof identifies the boundary maps as well as the central chain
groups. Without that step, shifting observed Betti numbers would be an
unjustified extrapolation. It also explains why the reference can be a
link of different component count: the argument is a statement about
marked complexes, not an isotopy between the original and reference
closures.

\begin{corollary}[Extraction from existing output]\label{cor:extraction}
The ordinary macro output at \(m_0\), which reports chain dimensions
and boundary ranks, determines \(\eta\) without a second homology
computation. Use
\[
h_*=\begin{cases}c+1&\varepsilon=+1,\\a-2&\varepsilon=-1,\end{cases}
\qquad
M=d_{h_*}=d_{h_*+1},\quad r=r_{h_*},\quad \eta=M-2r.
\]
\end{corollary}

The threshold Betti numbers at \(h_*\) and \(h_*+1\) need not equal
\(\eta\). They are boundary-region values. This distinction is covered
by tests comparing the extracted slope with the newly interior degree
at \(m_0+1\), in addition to tests of total ranks.

\subsection{Chain savings and a concrete family}

Let \(A_0=\sum_q\dim I^q\), \(B_0=\sum_q\dim E^q=M\).
The original and reference chain dimensions are exactly
\begin{equation}\label{eq:chain-saving}
\N(m)=A_0+mB_0,\qquad
\N_0=A_0+(W+2)B_0.
\end{equation}
The coefficient \(B_0\) counts repeated chain generators.
The coefficient \(\eta=B_0-2\rank F\) counts repeated homology.
They are usually different.

For example, take
\[
\beta_m=\sigma_1^m\sigma_2^{-1}\sigma_1\sigma_2^{-1}
\quad\text{on three strands},\qquad m\ge5.
\]
Here \((a,c,W,m_0)=(-2,1,3,5)\). The reference computation gives
\(\N_0=93\), \(B_0=15\), \(r=6\), and \(\eta=3\), while
\(\N(m)=15m+18\) and \(R(m)=3m+2\).
The result applies to every integer \(m\ge5\), including link parities;
odd \(m\) gives a knot. The figure included with the experiments plots
the full degree profile and its translated upper boundary.
This is a useful exact test family, not a hard recognition family:
the baseline already has structural and three-strand certificates.
